\chapter{Introducere}
\label{ch:intro}

\section{Context și motivație}

Testarea software a devenit una dintre activitățile esențiale în ingineria modernă a aplicațiilor. Costul unui defect descoperit târziu, în producție, poate fi de mai multe ori mai mare decât costul unui defect identificat în timpul dezvoltării \cite{myers2011}. Mai mult, în domeniile sensibile - financiar, medical, sau în servicii care manipulează date personale - bug-urile pot duce la pierderi materiale concrete și la deteriorarea încrederii utilizatorilor. Practic, software-ul fără teste e o promisiune care nu se poate verifica.

În contextul disciplinei \textit{Testarea Sistemelor Software}, proiectul nostru are ca scop aplicarea practică a strategiilor de testare unitară pe o clasă dintr-o aplicație reală. Nu pe un exemplu didactic simplificat, ci pe cod care funcționează efectiv într-o aplicație web și care prezintă toate caracteristicile incomode ale codului real: dependențe multiple, interacțiuni cu o bază de date, validări complexe, calcule monetare cu \texttt{BigDecimal}, condiții logice compuse cu trei sau patru operanzi.

Aplicația aleasă pentru studiul de caz se numește \textit{Călătorii} și e un planificator de vacanțe dezvoltat de noi pe stiva Spring Boot 3 + JPA + Thymeleaf. Are funcționalități multiple: gestiune călătorii, jurnal de călătorie, achiziție de tururi, leaderboard, hartă interactivă a țărilor vizitate, chatbot. Modulul pe care l-am ales pentru testare unitară este modulul \textit{Wallet}, mai precis clasa \texttt{WalletServiceImpl}, care răspunde de logica financiară a fiecărei călătorii.

\section{Obiectivele proiectului}

Cerințele proiectului T3 stabilesc, pentru fiecare echipă, următoarele obiective \cite{tss_assignment}:

\begin{itemize}
    \item Alegerea unei singure clase relevante și aplicarea unor teste unitare pe ea, ilustrând cele șase strategii fundamentale: partiționare în clase de echivalență, analiza valorilor de frontieră, acoperire la nivel de instrucțiune, decizie, condiție și pe circuite independente.
    \item Generarea unui raport de mutanți cu un instrument dedicat (\textsc{PITest} pentru Java) și omorârea explicită a cel puțin doi mutanți neechivalenți.
    \item Realizarea unui experiment cu un instrument AI care generează teste automat (ChatGPT, GitHub Copilot etc.), urmat de o comparație detaliată cu suita proprie.
    \item Documentarea completă a procesului, cu strategii, diagrame, configurația software și hardware, capturi de ecran cu rularea testelor, comparații tabelare și interpretări.
\end{itemize}

În plus, ne-am setat și câteva obiective interne, care depășesc strict cerințele:

\begin{itemize}
    \item Să atingem cel puțin 80\% pe \textit{branch coverage} măsurat cu JaCoCo, ca prag minim pentru o suită considerată solidă.
    \item Să atingem cel puțin 60\% mutation score, ținând cont de complexitatea ridicată a clasei alese.
    \item Să demonstrăm că aplicăm strategiile cu rigoare - nu doar din punct de vedere cantitativ, ci și conceptual, justificând în documentație de ce a fost ales fiecare test în parte.
\end{itemize}

\section{Echipa și diviziunea muncii}

Proiectul a fost realizat de o echipă de trei studenți. Munca a fost împărțită astfel încât fiecare membru să contribuie la fiecare etapă majoră:

\begin{itemize}
    \item \textbf{Iftime Raluca} (șef echipă) - dezvoltarea aplicației \textit{Călătorii}, integrarea modulului wallet, configurarea instrumentelor de testare (Maven, JaCoCo, PITest), redactarea documentației.
    \item \textbf{Lupeș Ioan-Marian} - dezvoltarea aplicației, integrarea funcționalităților secundare (chatbot, hartă), generarea capturilor de ecran și a videoclipurilor pentru livrabile.
    \item \textbf{Denisa Gheorghe} - dezvoltarea aplicației, contribuții la modulul wallet, suport pentru integrarea testelor și pentru experimentul cu AI.
\end{itemize}

\section{Structura lucrării}

Lucrarea este organizată în nouă capitole, urmate de o secțiune de bibliografie și o anexă cu fragmente de cod relevante. Pe scurt:

\begin{itemize}
    \item \textbf{Capitolul \ref{ch:preliminarii}} - prezentarea fundamentelor teoretice ale testării software, cu accent pe strategiile aplicate (Black Box / White Box / Mutation).
    \item \textbf{Capitolul \ref{ch:configurare}} - configurația de lucru: hardware, sistem de operare, versiunile instrumentelor folosite.
    \item \textbf{Capitolul \ref{ch:aplicatia}} - descrierea aplicației \textit{Călătorii} și a clasei \texttt{WalletServiceImpl}, cu motivarea alegerii.
    \item \textbf{Capitolul \ref{ch:strategii}} - aplicarea celor șase strategii pe clasa noastră, cu exemple, tabele, diagrame de control flow și fragmente de cod.
    \item \textbf{Capitolul \ref{ch:jacoco}} - rezultate JaCoCo (statement coverage și branch coverage), capturi de ecran, interpretarea procentelor neacoperite.
    \item \textbf{Capitolul \ref{ch:pitest}} - rezultate PITest, analiza mutanților supraviețuitori, mutanții neechivalenți pe care i-am omorât explicit.
    \item \textbf{Capitolul \ref{ch:ai}} - raport asupra folosirii unui tool AI (ChatGPT) pentru generarea automată de teste, cu prompt, răspuns, capturi de ecran și comparație cantitativă cu suita proprie.
    \item \textbf{Capitolul \ref{ch:concluzii}} - concluzii finale și direcții viitoare.
\end{itemize}

Anexa A conține fragmente de cod folosite ca referință pe parcursul documentației, iar bibliografia conține toate sursele consultate, inclusiv documentațiile oficiale ale instrumentelor folosite și articolele didactice de la curs.
